\documentclass[12pt, t, aspectratio=169]{beamer}
\usepackage{ctex}
\usepackage{amsmath, amssymb, amsthm}
\usepackage{graphicx}
\usepackage{booktabs}
\usepackage{listings}
\usepackage{xcolor}
\usepackage{etoolbox} % 用于条件判断
\usepackage{hyperref}
\usepackage{tikz}
\usetikzlibrary{positioning}

% 设置主题与配色
\usetheme{Madrid}
\usecolortheme{default}
\setbeamertemplate{navigation symbols}{} % 去除导航符号

% 标题页信息
\title{第二章：概率论与机器学习核心概念}
\subtitle{从贝叶斯定理到信息论}
% \institute{}
\author{《深度学习基础》课程小组}
\date{\today}

% 数学命令简写
\newcommand{\E}{\mathbb{E}}
\newcommand{\Var}{\text{Var}}
\newcommand{\Cov}{\text{Cov}}
\newcommand{\KL}{D_{\text{KL}}}

% 自定义命令：仅显示特定范围的 section 在目录中
% 用法: \tableofcontents[sections={1-10}]
\makeatletter
\newcommand{\tableofcontentsrange}[2]{%
  \begingroup
    \def\beamer@checksection##1{%
      \ifnum##1<#1\relax
        \let\beamer@next=\relax
      \else
        \ifnum##1>#2\relax
          \let\beamer@next=\relax
        \else
          \let\beamer@next=\beamer@showsection
        \fi
      \fi
      \beamer@next
    }%
    \tableofcontents
  \endgroup
}
\makeatother

\begin{document}

% 封面页
\begin{frame}
    \titlepage
\end{frame}

% 目录页 1 (第 1-10 个名词)
\begin{frame}{目录}
    \begin{columns}
        \column{0.5\textwidth}
        \tableofcontents[sections={1-10}]
        \column{0.5\textwidth}
        \tableofcontents[sections={11-20}]
    \end{columns}
%    \vspace{0.5cm}
%    \centering \small (续见下一页)
\end{frame}

% 目录页 2 (第 11-20 个名词)
%\begin{frame}{目录 (2/2)}
%    \begin{columns}
%        \column{0.5\textwidth}
%        \tableofcontents[sections={11-15}]
%        \column{0.5\textwidth}
%        \tableofcontents[sections={16-20}]
%    \end{columns}
%\end{frame}

% --- 以下内容保持不变 ---

% 1. Bayes' theorem
\section{Bayes' theorem}
\begin{frame}{贝叶斯定理 (Bayes' Theorem)}
    \begin{block}{定义}
        贝叶斯定理描述了在已知相关证据 $B$ 的情况下，事件 $A$ 发生的概率如何更新。它是贝叶斯统计推断的核心。
        它提供了一种根据新数据更新信念（概率）的数学框架。
    \end{block}
    
    \begin{equation*}
        P(A|B) = \frac{P(B|A)P(A)}{P(B)}
    \end{equation*}
    
    % \vspace{0.5cm}
    \begin{itemize}
        \item $P(A|B)$: 后验概率 (Posterior)
        \item $P(B|A)$: 似然函数 (Likelihood)
        \item $P(A)$: 先验概率 (Prior)
        \item $P(B)$: 证据或边际概率 (Evidence)
        % \item 意义：
    \end{itemize}
    
\end{frame}

% 2. prior probability
\section{prior probability}
\begin{frame}{先验概率 (Prior Probability)}
    \begin{block}{定义}
        在观察到任何新数据或证据之前，基于现有知识、历史数据或主观信念对事件发生概率的估计。
    \end{block}
    
    %\vspace{0.5cm}
    \begin{itemize}
        \item 记作 $P(\theta)$ 或 $P(A)$。
        \item 代表“事前”的信念。
        \item 在贝叶斯推断中，先验可以是无信息的（Uniform），也可以是有信息的（Informative）。
    \end{itemize}
    
    % \vspace{0.5cm}
    \begin{exampleblock}{示例}
        在抛硬币实验中，如果我们认为硬币是公平的，那么正面朝上的先验概率 $P(H) = 0.5$。
    \end{exampleblock}
\end{frame}

% 3. posterior probability
\section{posterior probability}
\begin{frame}{后验概率 (Posterior Probability)}
    \begin{block}{定义}
        在考虑了新的证据或数据之后，对事件发生概率的修正估计。
    \end{block}
    
    \vspace{0.2cm}
    \begin{itemize}
        \item 记作 $P(\theta|D)$ 或 $P(A|B)$。
        \item 是贝叶斯推断的最终目标。
        \item 计算公式：$\text{后验} \propto \text{似然} \times \text{先验}$。
    \end{itemize}
    
    \vspace{0.2cm}
    \begin{alertblock}{直观理解}
        后验概率结合了我们的初始信念（先验）和新观察到的数据（似然），反映了更新后的认知状态。
    \end{alertblock}
\end{frame}

% 4. Independent variables
\section{Independent variables}
\begin{frame}{独立变量 (Independent Variables)}
    \begin{block}{定义}
        如果两个随机变量 $X$ 和 $Y$ 的联合概率分布等于它们边缘概率分布的乘积，则称它们是统计独立的。
    \end{block}
    
    \begin{equation*}
        P(X, Y) = P(X)P(Y)
    \end{equation*}
    或者条件概率形式：
    \begin{equation*}
        P(X|Y) = P(X) \quad (\text{当 } P(Y) > 0)
    \end{equation*}
    
    \vspace{0.5cm}
    \begin{itemize}
        \item 独立性意味着知道 $Y$ 的值不会提供关于 $X$ 的任何信息。
        \item 注意：独立一定不相关，但不相关不一定独立（除非是高斯分布）。
    \end{itemize}
\end{frame}

% 5. probability density function
\section{probability density function}
\begin{frame}{概率密度函数 (PDF)}
    \begin{block}{定义}
        对于连续型随机变量 $X$，概率密度函数 $f(x)$ 描述了变量在某个特定值附近的相对可能性。
    \end{block}
    
    \begin{itemize}
        \item $f(x) \geq 0$
        \item $\int_{-\infty}^{\infty} f(x) dx = 1$
        \item 变量落在区间 $[a, b]$ 的概率为：
        %\begin{equation*}
         $\displaystyle   P(a \leq X \leq b) = \int_{a}^{b} f(x) dx$
        %\end{equation*}
    \end{itemize}
    
    \vspace{0.2cm}
    \begin{alertblock}{注意}
        $f(x)$ 本身不是概率（可以大于1），只有积分面积才是概率。
    \end{alertblock}
\end{frame}

% 6. cumulative distribution function
\section{cumulative distribution function}
\begin{frame}{累积分布函数 (CDF)}
    \begin{block}{定义}
        累积分布函数 $F(x)$ 表示随机变量 $X$ 取值小于或等于 $x$ 的概率。
    \end{block}
    
    \begin{equation*}
        F(x) = P(X \leq x)
    \end{equation*}
    
    \vspace{0.5cm}
    \begin{itemize}
        \item 对于连续变量：$F(x) = \int_{-\infty}^{x} f(t) dt$。
        \item 性质：单调不减，$F(-\infty)=0$, $F(\infty)=1$。
        \item PDF 是 CDF 的导数：$f(x) = \frac{d}{dx}F(x)$。
    \end{itemize}
\end{frame}

% 7. expectation
\section{expectation}
\begin{frame}{期望 (Expectation)}
    \begin{block}{定义}
        期望（或均值）是随机变量取值的加权平均，权重为对应的概率。它描述了分布的中心位置。
    \end{block}
    
    \begin{itemize}
        \item 离散型：$\E[X] = \sum_{x} x P(X=x)$
        \item 连续型：$\E[X] = \int_{-\infty}^{\infty} x f(x) dx$
        \item 线性性质：$\E[aX + bY] = a\E[X] + b\E[Y]$
    \end{itemize}
    
    \vspace{0.2cm}
    \begin{exampleblock}{意义}
        如果在大量重复实验中观测随机变量，其算术平均值将收敛于期望值（大数定律）。
    \end{exampleblock}
\end{frame}

% 8. covariance
\section{covariance}
\begin{frame}{协方差 (Covariance)}
    \begin{block}{定义}
        协方差衡量两个随机变量 $X$ 和 $Y$ 之间的线性相关程度及其方向。
    \end{block}
    
    \begin{equation*}
        \Cov(X, Y) = \E[(X - \E[X])(Y - \E[Y])] = \E[XY] - \E[X]\E[Y]
    \end{equation*}
    
    \vspace{0.5cm}
    \begin{itemize}
        \item $\Cov(X, Y) > 0$: 正相关（同向变化）。
        \item $\Cov(X, Y) < 0$: 负相关（反向变化）。
        \item $\Cov(X, Y) = 0$: 不相关（但不一定独立）。
        \item $\Var(X) = \Cov(X, X)$。
    \end{itemize}
\end{frame}

% 9. conditional expectation
\section{conditional expectation}
\begin{frame}{条件期望 (Conditional Expectation)}
    \begin{block}{定义}
        在已知另一个随机变量 $Y$ 取特定值 $y$ 的条件下，随机变量 $X$ 的期望值。
    \end{block}
    
    \begin{equation*}
        \E[X|Y=y] = \begin{cases} 
        \sum_{x} x P(X=x|Y=y) & \text{(离散)} \\
        \int_{-\infty}^{\infty} x f_{X|Y}(x|y) dx & \text{(连续)}
        \end{cases}
    \end{equation*}
    
    \vspace{0.5cm}
    \begin{alertblock}{全期望公式}
        $\E[X] = \E[\E[X|Y]]$。即：无条件期望等于条件期望的期望。
    \end{alertblock}
\end{frame}

% 10. Gaussian distribution
\section{Gaussian distribution}
\begin{frame}{高斯分布 (Gaussian Distribution)}
    \begin{block}{定义}
        又称正态分布，是最常见的连续概率分布，由均值 $\mu$ 和方差 $\sigma^2$ 完全确定。
    \end{block}
    
    \begin{equation*}
        \mathcal{N}(x|\mu, \sigma^2) = \frac{1}{\sqrt{2\pi\sigma^2}} \exp\left(-\frac{(x-\mu)^2}{2\sigma^2}\right)
    \end{equation*}
    
    \vspace{0.5cm}
    \begin{itemize}
        \item 钟形曲线，关于 $\mu$ 对称。
        \item 中心极限定理：大量独立随机变量之和趋向于高斯分布。
        \item 在机器学习中广泛用于噪声建模和先验假设。
    \end{itemize}
\end{frame}

% 11. likelihood function
\section{likelihood function}
\begin{frame}{似然函数 (Likelihood Function)}
    \begin{block}{定义}
        给定观测数据 $D$ 时，关于模型参数 $\theta$ 的函数。它衡量参数 $\theta$ 产生该数据的可能性。
    \end{block}
    
    \begin{equation*}
        L(\theta|D) = P(D|\theta)
    \end{equation*}
    
    \vspace{0.1cm}
    \begin{itemize}
        \item 注意：虽然形式上等于条件概率，但视为 $\theta$ 的函数时，它不是概率分布（对 $\theta$ 积分不一定为1）。
        \item \textbf{最大似然估计 (MLE)}: 寻找使 $L(\theta|D)$ 最大的 $\theta$。
    \end{itemize}
    
    \begin{alertblock}{区别}
        概率是固定参数预测数据；似然是固定数据评估参数。
    \end{alertblock}
\end{frame}

% 12. problem of over-fitting
\section{problem of over-fitting}
\begin{frame}{过拟合问题 (Problem of Over-fitting)}
    \begin{block}{定义}
        模型在训练数据上表现极好（误差很低），但在未见过的测试数据上表现很差的现象。
    \end{block}
    
    \vspace{0.5cm}
    \begin{itemize}
        \item \textbf{原因}: 模型过于复杂，记住了训练数据中的噪声而非潜在规律。
        \item \textbf{特征}: 高方差 (High Variance)，低偏差 (Low Bias)。
        \item \textbf{解决方法}:
        \begin{itemize}
            \item 正则化 (Regularization, e.g., L1, L2)
            \item 增加训练数据
            \item 降低模型复杂度
            \item Dropout (神经网络)
            \item 交叉验证
        \end{itemize}
    \end{itemize}
\end{frame}

% 13. sum-of-squares error function
\section{sum-of-squares error function}
\begin{frame}{平方和误差函数 (Sum-of-Squares Error)}
    \begin{block}{定义}
        常用于回归问题的损失函数，衡量模型预测值 $y(x, w)$ 与真实目标值 $t$ 之间的差异。
    \end{block}
    
    \begin{equation*}
        E(w) = \frac{1}{2} \sum_{n=1}^{N} \{ y(x_n, w) - t_n \}^2
    \end{equation*}
    
    \vspace{0.5cm}
    \begin{itemize}
        \item 假设噪声服从高斯分布时，最小化 SSE 等价于最大似然估计。
        \item 系数 $\frac{1}{2}$ 是为了求导方便。
        \item 对异常值（Outliers）敏感。
    \end{itemize}
\end{frame}

% 14. entropy of the random variable
\section{entropy of the random variable}
\begin{frame}{随机变量的熵 (Entropy)}
    \begin{block}{定义}
        熵 $H(X)$ 衡量离散随机变量 $X$ 的不确定性或信息含量。
    \end{block}
    
    \begin{equation*}
        H(X) = - \sum_{x} P(x) \log_2 P(x)
    \end{equation*}
    
    \vspace{0.2cm}
    \begin{itemize}
        \item 单位通常为比特 (bits)。
        \item $H(X) \geq 0$。
        \item 当分布均匀时熵最大（最不确定）；当某事件概率为1时熵为0（确定）。
        \item 表示编码该变量所需的最小平均比特数。
    \end{itemize}
\end{frame}

% 15. differential entropy
\section{differential entropy}
\begin{frame}{微分熵 (Differential Entropy)}
    \begin{block}{定义}
        连续随机变量 $X$ 的熵的推广，称为微分熵 $h(X)$。
    \end{block}
    
    \begin{equation*}
        h(X) = - \int_{-\infty}^{\infty} f(x) \log f(x) dx
    \end{equation*}
    
    \vspace{0.2cm}
    \begin{alertblock}{注意}
        \begin{itemize}
            \item 微分熵可以是负数（不同于离散熵）。
            \item 它不是绝对信息量的度量，而是相对于均匀分布的相对度量。
            \item 坐标变换下不具有不变性。
        \end{itemize}
    \end{alertblock}
\end{frame}

% 16. relative entropy or Kullback–Leibler divergence
\section{relative entropy or Kullback–Leibler divergence}
\begin{frame}{相对熵 / KL 散度 (KL Divergence)}
    \begin{block}{定义}
        衡量两个概率分布 $P$ (真实分布) 和 $Q$ (近似分布) 之间差异的非对称度量。
    \end{block}
    
    \begin{equation*}
        \KL(P||Q) = \sum_{x} P(x) \log \frac{P(x)}{Q(x)} = \E_{P} \left[ \log \frac{P(x)}{Q(x)} \right]
    \end{equation*}
    
    \vspace{0.2cm}
    \begin{itemize}
        \item $\KL(P||Q) \geq 0$，当且仅当 $P=Q$ 时取等号。
        \item 非对称：$\KL(P||Q) \neq \KL(Q||P)$。
        \item 广泛应用于变分推断、生成模型（如 VAE）和强化学习。
    \end{itemize}
\end{frame}

% 17. conditional entropy
\section{conditional entropy}
\begin{frame}{条件熵 (Conditional Entropy)}
    \begin{block}{定义}
        在已知随机变量 $Y$ 的条件下，随机变量 $X$ 剩余的不确定性。
    \end{block}
    
    \begin{equation*}
        H(X|Y) = \sum_{y} P(y) H(X|Y=y) = - \sum_{x,y} P(x,y) \log P(x|y)
    \end{equation*}
    
    \vspace{0.2cm}
    \begin{itemize}
        \item 链式法则：$H(X, Y) = H(Y) + H(X|Y)$。
        \item $H(X|Y) \leq H(X)$：知道 $Y$ 的信息永远不会增加 $X$ 的不确定性。
        \item 若 $X, Y$ 独立，则 $H(X|Y) = H(X)$。
    \end{itemize}
\end{frame}

% 18. mutual information
\section{mutual information}
\begin{frame}{互信息 (Mutual Information)}
    \begin{block}{定义}
        衡量两个随机变量 $X$ 和 $Y$ 之间相互依赖程度的量。表示知道其中一个变量能减少另一个变量多少不确定性。
    \end{block}
    
    \begin{equation*}
        I(X; Y) = H(X) - H(X|Y) = H(Y) - H(Y|X)
    \end{equation*}
    或者用 KL 散度表示：
    \begin{equation*}
        I(X; Y) = \KL(P(X,Y) || P(X)P(Y))
    \end{equation*}
    
    \vspace{0.2cm}
    \begin{itemize}
        \item $I(X; Y) \geq 0$。
        \item 若 $X, Y$ 独立，则 $I(X; Y) = 0$。
        \item 对称性：$I(X; Y) = I(Y; X)$。
    \end{itemize}
\end{frame}

% 19. maximum a posteriori estimate
\section{maximum a posteriori estimate}
\begin{frame}{最大后验估计 (MAP Estimate)}
    \begin{block}{定义}
        一种参数估计方法，通过最大化后验概率来寻找最优参数 $\theta$。
    \end{block}
    
    \begin{equation*}
        \hat{\theta}_{\text{MAP}} = \arg\max_{\theta} P(\theta|D) = \arg\max_{\theta} P(D|\theta)P(\theta)
    \end{equation*}
    
    \vspace{0.2cm}
    \begin{itemize}
        \item 结合了似然 $P(D|\theta)$ 和先验 $P(\theta)$。
        \item 是贝叶斯学派的一种点估计方法。
        \item 当先验为均匀分布时，MAP 退化为最大似然估计 (MLE)。
        \item 正则化项通常对应于先验分布（例如 L2 正则化对应高斯先验）。
    \end{itemize}
\end{frame}

% 20. Bayesian machine learning
\section{Bayesian machine learning}
\begin{frame}{贝叶斯机器学习 (Bayesian Machine Learning)}
    \begin{block}{定义}
        一种机器学习方法论，将模型参数视为随机变量，利用贝叶斯定理结合先验知识和观测数据来推断参数的后验分布，而非单一的点估计。
    \end{block}
    
    \vspace{0.2cm}
    \begin{itemize}
        \item \textbf{核心优势}:
        \begin{itemize}
            \item 自然地处理不确定性（Uncertainty Quantification）。
            \item 有效防止过拟合（通过积分边缘化参数）。
            \item 可以融入领域知识（通过先验）。
        \end{itemize}
        \item \textbf{挑战}: 后验分布的计算通常难以解析求解，需借助近似推断（如 MCMC, 变分推断 VI）。
    \end{itemize}
\end{frame}

\end{document}

做一个latex beamer, 解释深度学习基础与概念里的有关概率理论的关键名词，每个名词一页slide（必要时可多页）， 请输出中文 latex 代码。以下是部分术语，请参考：

Bayes’ theorem

prior probability

posterior probability

Independent variables

probability density function

cumulative distribution function

expectation

covariance

conditional expectation

Gaussian distribution

likelihood function

problem of over-fitting

sum-of-squares error function

entropy of the random variable

differential entropy

relative entropy or Kullback–Leibler divergence

conditional entropy

mutual information

maximum a posteriori estimate

Bayesian machine learning


